\chapter{CONCLUSION AND FUTURE WORK}

%-------------------------------------------------------------

\section{Conclusion}

In this project, a Machine Learning-based message 
prioritization and congestion-aware rate control 
system for UAV communication networks was proposed 
and implemented. The main objective of the project 
was to improve communication reliability and reduce 
packet loss in MAVLink-based UAV networks operating 
under dynamic network conditions.

The proposed system integrates Random Forest-based 
message classification with adaptive rate control 
mechanisms. The Random Forest classifier was used 
to assign priority levels to MAVLink messages based 
on their importance and communication requirements. 
High-priority messages such as emergency commands 
and navigation updates were identified and given 
transmission preference over low-priority messages.

The final version of the project also extends the 
communication pipeline with reinforcement learning-based 
scheduling and packet aggregation. In this combined 
framework, the learned policy jointly decides which 
priority queue to serve and how aggressively to aggregate 
packets before transmission. This makes the overall system 
more adaptive than a fixed rule-based scheduler.

The congestion-aware rate control mechanism was 
designed to monitor network parameters such as 
queue occupancy, packet loss ratio, and bandwidth 
utilization. Based on congestion levels, the system 
dynamically adjusted transmission rates to reduce 
network overload. Priority-based queue scheduling 
was used to ensure timely delivery of critical 
messages.

The closed-loop execution structure created for the 
combined model links the policy layer and the simulation 
layer in a feedback cycle: the current state is observed, 
an action is selected, the action is applied, and the next 
state is generated from the resulting traffic condition. 
This loop is what allows the project to evaluate different 
trained checkpoints against the same baseline conditions.

The implementation of the proposed system was 
successfully carried out using Python programming 
and machine learning libraries. Simulation results 
demonstrated that the proposed approach significantly 
improves network performance compared to traditional 
fixed-rate communication systems.

The major achievements of the project include:

\begin{itemize}

\item Successful implementation of Random Forest-based 
message prioritization.

\item Effective congestion detection and adaptive 
rate control.

\item Joint scheduling and aggregation using a learned 
reinforcement learning policy.

\item Reduction in packet loss under varying 
traffic conditions.

\item Improvement in throughput and communication 
reliability while staying close to the paper baseline 
in the best-performing RL checkpoint.

\item Reduction in transmission delay for 
high-priority messages.

\end{itemize}

Overall, the proposed system enhances the Quality 
of Service (QoS) of UAV communication networks 
and improves the reliability of MAVLink-based 
message transmission.


%-------------------------------------------------------------

\section{Future Work}

Although the proposed system provides significant 
improvements in communication efficiency, there 
are several areas where further enhancements can 
be made in future research.

Possible future improvements include:

\begin{itemize}

\item Implementation of the proposed system in 
real UAV hardware environments to validate 
performance under real-world conditions.

\item Extension of the closed-loop RL policy to 
multi-UAV cooperative scheduling and aggregation.

\item Integration of advanced Machine Learning 
techniques such as Deep Learning models for 
more accurate message classification.

\item Development of energy-efficient communication 
strategies to reduce power consumption in UAV 
networks.

\item Extension of the system to support large-scale 
multi-UAV networks with higher node density.

\item Incorporation of security mechanisms to 
protect communication from cyber threats and 
unauthorized access.

\item Implementation of adaptive routing protocols 
combined with congestion control for improved 
network scalability.

\item Online adaptation of the learned policy so that 
the scheduler can refresh its decisions when traffic 
patterns drift over time.

\end{itemize}

Future developments in intelligent communication 
systems and Machine Learning techniques are 
expected to further enhance UAV communication 
efficiency and reliability.

%-------------------------------------------------------------

\section{Summary}

This chapter presented the overall conclusion 
of the project and highlighted the key outcomes 
achieved through the implementation of the 
proposed system. The future work section 
outlined possible directions for further 
improvements and research extensions in 
UAV communication networks.